\chapter{Behçet's Disease}
\index{Behçet's Disease}
\label{sec:Behçet's Disease}

\section{Overview}
\begin{itemize}[noitemsep]
  \item Behçet's disease is a chronic, relapsing inflammatory disorder causing inflammation of blood vessels (vasculitis) throughout the body
  \item It classically presents with a triad of recurrent mouth ulcers, genital ulcers, and eye inflammation (uveitis)
  \item It is most common along the historic "Silk Road" (Turkey, Middle East, East Asia) and can affect any organ system
\end{itemize}
\section{Causes}
The exact cause is unknown, but the disease is thought to involve an abnormal immune response triggered by infection or other environmental factors in genetically susceptible people. The HLA-B51 gene is strongly associated with the disease. The immune system attacks the body's own blood vessels, causing inflammation, thrombosis, and tissue damage.
\section{Risk Factors}
\begin{itemize}[noitemsep]
  \item Ethnic background (Turkish, Middle Eastern, East Asian ancestry)
  \item HLA-B51 gene carriage
  \item Male sex (associated with more severe eye disease)
  \item Young adulthood, typically ages 20--40
  \item Family history of the disease
\end{itemize}
\section{Signs and Symptoms}
\subsection{Common symptoms}
\begin{itemize}[noitemsep]
  \item Recurrent painful mouth ulcers (aphthous ulcers)
  \item Painful genital ulcers
  \item Eye inflammation: uveitis, redness, pain, blurred vision, floaters
  \item Skin lesions, including acne-like bumps and nodules
  \item Joint pain and swelling (arthritis)
  \item Fatigue and fever
\end{itemize}
\subsection{Emergency warning signs}
\begin{itemize}[noitemsep]
  \item Sudden vision loss or severe eye pain requires urgent evaluation, as eye involvement can cause blindness
  \item Sudden headache, weakness, or confusion may indicate nervous system involvement
  \item Chest pain, shortness of breath, or coughing up blood may indicate vascular involvement
\end{itemize}
\section{Progression and Stages}
The disease follows a relapsing-remitting course with unpredictable flares and periods of improvement. It tends to be most severe in young men. Ocular disease, if untreated, causes progressive inflammation that can lead to scarring, glaucoma, and blindness. Vascular complications such as venous thrombosis and arterial aneurysm can occur and are potentially life-threatening.
\section{Complications}
\begin{itemize}[noitemsep]
  \item Permanent vision loss and blindness
  \item Central nervous system involvement (meningoencephalitis, stroke)
  \item Venous thrombosis, including deep vein and cerebral venous thrombosis
  \item Aneurysm or rupture of major arteries
  \item Gastrointestinal ulceration and bleeding
  \item Chronic joint damage and disability
\end{itemize}
\section{Diagnosis}
\begin{itemize}[noitemsep]
  \item Clinical diagnosis based on the characteristic pattern of symptoms (recurrent oral ulcers plus at least two of: genital ulcers, eye inflammation, skin lesions, or positive pathergy test)
  \item Pathergy test (skin response to a small needle prick)
  \item HLA-B51 genetic testing as a supporting clue
  \item Eye examination with slit lamp and fundoscopy
  \item Blood tests for inflammation (ESR, CRP) and imaging as indicated
  \item Exclusion of other causes of uveitis and vasculitis
\end{itemize}
\section{Prevention}
\begin{itemize}[noitemsep]
  \item There is no known prevention for Behçet's disease
  \item Early recognition and treatment reduce the risk of irreversible organ damage
  \item Regular ophthalmological follow-up to detect silent eye inflammation
  \item Managing cardiovascular risk factors
\end{itemize}
\section{Treatment Options}
\begin{itemize}[noitemsep]
  \item Topical or oral corticosteroids for mouth and genital ulcers
  \item Colchicine for mucosal and joint symptoms
  \item Immunosuppressants (azathioprine, cyclosporine, methotrexate) for moderate to severe disease
  \item Biologic agents (infliximab, adalimumab, secukinumab, interferon) for refractory or sight-threatening disease
  \item Anticoagulation for thrombosis (with caution) alongside immunosuppression
  \item Regular multidisciplinary care (ophthalmology, rheumatology, neurology)
\end{itemize}
\section{Living with It}
\begin{itemize}[noitemsep]
  \item Follow the treatment plan and attend regular specialist follow-up
  \item Protect vision with regular eye examinations and prompt reporting of symptoms
  \item Manage flares with the help of the care team
  \item Seek support from patient organizations
  \item Maintain a healthy lifestyle and manage stress
\end{itemize}
\section{Outlook}
With modern immunosuppressive and biologic therapy, the outlook has improved substantially. Many patients achieve good control, but the disease is lifelong and relapsing. Eye involvement remains the most serious complication; early aggressive treatment significantly reduces the risk of blindness.
